\documentclass[10pt,a4paper]{amsart}
\usepackage[margin=19mm]{geometry}
\usepackage{amsmath,amssymb,mathtools}
\usepackage[scr=rsfs,cal=euler]{mathalfa}
\usepackage{xfrac}
\usepackage[unicode,hidelinks,bookmarks=false]{hyperref}
\newcommand{\ie}{i.e.\ }
\newcommand{\cf}{cf.\ }
\newcommand{\resp}{respectively\ }
\newcommand{\Subsection}{\subsection}
\providecommand{\nobreakdash}{\nobreak\hbox{-}}
\setlength{\emergencystretch}{2em}
\multlinegap0pt
\usepackage{amssymb}
\usepackage{mathtools}
\newtheorem{theorem}{Theorem}
\newcommand{\IS}{\operatorname{IS}}
\newcommand{\lub}{\operatorname{lub}}
\newcommand{\succop}{\operatorname{succ}}
\title{Bourbaki--Zorn Normal Forms for Maximality Arguments}
\author{You-Chang Liu}
\date{Source: arXiv:2605.08274v2 (CC BY 4.0)}
\begin{document}
\maketitle

\noindent\emph{Source and licence.} Bourbaki--Zorn Normal Forms for Maximality Arguments. Version arXiv:2605.08274v2 (CC BY 4.0), distributed by arXiv under the Creative Commons Attribution 4.0 International licence, http://creativecommons.org/licenses/by/4.0/, as recorded in the arXiv licence metadata for this exact version and shown on the abstract page https://arxiv.org/abs/2605.08274v2. Full licence notice: \texttt{licenses/arxiv-2605.08274-CC-BY-4.0.txt}.\par\smallskip

\noindent\textit{Editorial scope.} Counted unit: the article's theorem theorem (source0.tex lines 401--423). The article's complete original proof of it is delivered with it (source0.tex lines 425--592, one \texttt{proof} environment of 168 lines). No mathematics is written, repaired or re-derived: the statement and the proof are the article's own lines, in the article's own notation, and no retained line is substituted or re-flowed.  No further source-local context is needed: every cross-reference of every retained line resolves inside the two delivered spans.  Nothing was omitted from any retained span, so the package carries no omission record.  No retained line contains a citation, so the wrapper carries no bibliography and no bibliography entry.  No retained line contains a figure environment, an image-inclusion command or drawing code, so no image is delivered and the package declares no asset dependency.  Editorial formatting only, in addition: an AMS \texttt{amsart} preamble that restates the article's own declarations and macros used by the retained lines, namely the environments \texttt{theorem} on separate counters, together with the standard packages those lines need.  The retained environments print in this document as Theorem 1 in order of appearance, whereas the article prints the same items with the numbering of its own full document.  The complete native source of the article as distributed in the arXiv e-print for this exact version is bundled unchanged as \texttt{sources/200144/source0.tex}.
\begin{theorem}[Well-ordered Bourbaki--Witt normal form]
Let \((X,\leq)\) be a nonempty partially ordered set such that every nonempty
well-ordered subset of \(X\) admits a least upper bound in \(X\). Let \(f:X\to X\) be
progressive, and fix \(x_0\in X\).

Then there exists a largest Bourbaki \(f\)-tower based at \(x_0\), denoted
\[
\Omega_f(x_0).
\]
Moreover, if
\[
\omega=\lub_X(\Omega_f(x_0)),
\]
then
\[
\omega\in \Omega_f(x_0)
\]
and
\[
f(\omega)=\omega.
\]
In particular, every progressive self-map on \(X\) has a fixed point.
\end{theorem}

\begin{proof}
Let \(\mathcal T\) be the collection of all Bourbaki \(f\)-towers based at \(x_0\). The
singleton set \(\{x_0\}\) belongs to \(\mathcal T\), so \(\mathcal T\neq\varnothing\).

By the Tower Comparison Lemma, any two elements of \(\mathcal T\) are comparable by
the initial-segment relation. Define
\[
\Omega_f(x_0)=\bigcup_{Y\in\mathcal T}Y.
\]

We verify the required properties in stages.

\begin{enumerate}
\item \emph{\(\Omega_f(x_0)\) is well-ordered.}

Since any two members of \(\mathcal T\) are comparable by the initial-segment relation,
the orders on the towers are compatible. Hence the order induced from \(X\) is total on
\(\Omega_f(x_0)\).

Let \(A\subseteq \Omega_f(x_0)\) be nonempty. Choose \(a\in A\), and choose
\(Y_a\in\mathcal T\) such that \(a\in Y_a\). Since \(A\cap Y_a\) is a nonempty subset of
the well-ordered set \(Y_a\), it has a least element; call it \(a_0\).

We claim that \(a_0\) is the least element of \(A\). Let \(b\in A\). If \(b\in Y_a\), then
\(a_0\leq b\) by the definition of \(a_0\). If \(b\notin Y_a\), choose
\(Y_b\in\mathcal T\) such that \(b\in Y_b\). By the Tower Comparison Lemma, \(Y_a\) and
\(Y_b\) are comparable by the initial-segment relation. Since \(b\notin Y_a\), the only
possible case is that \(Y_a\) is a proper initial segment of \(Y_b\). Hence every element
of \(Y_a\) is strictly below \(b\). In particular,
\[
a_0<b.
\]
Thus \(a_0\leq b\) for every \(b\in A\), so \(a_0\) is the least element of \(A\).
Therefore \(\Omega_f(x_0)\) is well-ordered. Its least element is \(x_0\).

\item \emph{\(\Omega_f(x_0)\) satisfies the successor condition.}

Let \(y\in\Omega_f(x_0)\), and suppose that \(y\) is not the largest element of
\(\Omega_f(x_0)\). Then there exists \(z\in\Omega_f(x_0)\) such that \(y<z\). Choose
towers \(Y,Z\in\mathcal T\) such that \(y\in Y\) and \(z\in Z\). By the Tower Comparison
Lemma, \(Y\) and \(Z\) are comparable by the initial-segment relation. Since \(y<z\), one
of these two towers contains both \(y\) and \(z\). In that tower, \(y\) is not largest, and
its successor is \(f(y)\). Hence
\[
f(y)\in\Omega_f(x_0).
\]

Now let \(u\in\Omega_f(x_0)\) satisfy \(y<u\). By the same comparability argument, there
is a tower in \(\mathcal T\) containing both \(y\) and \(u\). In that tower, the least
element strictly above \(y\) is \(f(y)\). Therefore
\[
f(y)\leq u.
\]
Thus \(f(y)\) is the least element of \(\Omega_f(x_0)\) strictly above \(y\). Hence
\[
\succop_{\Omega_f(x_0)}(y)=f(y).
\]

\item \emph{\(\Omega_f(x_0)\) satisfies the limit condition.}

Let \(y\) be a limit element of \(\Omega_f(x_0)\). Choose \(Y\in\mathcal T\) such that
\(y\in Y\). We claim that \(y\) is a limit element of \(Y\).

First, \(y\) cannot be the least element of \(Y\), because then \(y=x_0\), which is also
the least element of \(\Omega_f(x_0)\), contrary to the assumption that \(y\) is a limit
element of \(\Omega_f(x_0)\). Next, suppose that \(y\) is a successor element of \(Y\),
say
\[
y=\succop_Y(u).
\]
We show that \(y\) would then be a successor element of \(\Omega_f(x_0)\), giving a
contradiction. Let \(v\in\Omega_f(x_0)\) satisfy
\[
u<v<y.
\]
Choose \(Z\in\mathcal T\) such that \(v\in Z\). By the Tower Comparison Lemma, \(Y\)
and \(Z\) are comparable by the initial-segment relation. If \(Z\) is an initial segment of
\(Y\), then \(v\in Y\), contradicting the fact that \(y\) is the successor of \(u\) in
\(Y\). If \(Y\) is an initial segment of \(Z\), then \(u,y\in Z\), and since \(v<y\), the
initial-segment property forces \(v\in Y\), again a contradiction. Thus there is no element
of \(\Omega_f(x_0)\) strictly between \(u\) and \(y\). Hence \(y\) is the successor of
\(u\) in \(\Omega_f(x_0)\), contradicting the assumption that \(y\) is a limit element of
\(\Omega_f(x_0)\). Therefore \(y\) is a limit element of \(Y\).

We now compare initial segments. If \(v\in\IS_{\Omega_f(x_0)}(y)\), choose
\(Z\in\mathcal T\) such that \(v\in Z\). By the Tower Comparison Lemma, \(Y\) and \(Z\)
are comparable by the initial-segment relation. If \(Z\) is an initial segment of \(Y\),
then \(v\in Y\). If \(Y\) is an initial segment of \(Z\), then \(y\in Y\subseteq Z\), and
\(v<y\) implies \(v\in Y\). Hence \(v\in\IS_Y(y)\). The reverse inclusion is immediate.
Therefore
\[
\IS_{\Omega_f(x_0)}(y)=\IS_Y(y).
\]
Since \(Y\) is a Bourbaki tower and \(y\) is a limit element of \(Y\),
\[
y=\lub_X(\IS_Y(y)).
\]
Consequently,
\[
y=\lub_X(\IS_{\Omega_f(x_0)}(y)).
\]

\item \emph{\(\Omega_f(x_0)\) is the largest Bourbaki tower.}

By the preceding three steps, \(\Omega_f(x_0)\) is a Bourbaki \(f\)-tower. Therefore
\[
\Omega_f(x_0)\in\mathcal T.
\]
It is clearly the largest element of \(\mathcal T\) under inclusion.

\item \emph{The least upper bound of \(\Omega_f(x_0)\) belongs to the tower.}

Since \(\Omega_f(x_0)\) is a nonempty well-ordered subset of \(X\), the hypothesis gives
a least upper bound in \(X\). Let
\[
\omega=\lub_X(\Omega_f(x_0)).
\]

Suppose first that \(\omega\notin\Omega_f(x_0)\). Since \(\omega\) is the least upper
bound of \(\Omega_f(x_0)\), every element of \(\Omega_f(x_0)\) is strictly below
\(\omega\). Moreover, \(\Omega_f(x_0)\) cannot have a largest element; otherwise its
largest element would already be the least upper bound of \(\Omega_f(x_0)\), contradicting
\(\omega\notin\Omega_f(x_0)\).

Adjoin \(\omega\) to \(\Omega_f(x_0)\) as a new largest element. The resulting set is
well-ordered. Since \(\Omega_f(x_0)\) has no largest element, the newly adjoined element
\(\omega\) is a limit element of the enlarged tower, and
\[
\omega=\lub_X(\Omega_f(x_0)).
\]
All earlier successor and limit conditions are unchanged. Hence the enlarged set is again
a Bourbaki \(f\)-tower based at \(x_0\), contradicting the maximality of
\(\Omega_f(x_0)\). Therefore
\[
\omega\in\Omega_f(x_0).
\]

\item \emph{The element \(\omega\) is a fixed point of \(f\).}

Suppose that
\[
f(\omega)>\omega.
\]
Since \(\omega\) is an upper bound of \(\Omega_f(x_0)\), every element of
\(\Omega_f(x_0)\) is below \(\omega\), and hence strictly below \(f(\omega)\). Therefore
\[
\Omega_f(x_0)\cup\{f(\omega)\}
\]
is obtained by adjoining \(f(\omega)\) as a new largest element.

In the enlarged set, \(f(\omega)\) is the successor of \(\omega\). The successor condition
therefore holds at \(\omega\), since its new successor is precisely \(f(\omega)\). The
newly added element \(f(\omega)\) is a successor element, so no new limit condition has to
be verified at that stage. All earlier successor and limit conditions are unchanged. Hence
the enlarged set is again a Bourbaki \(f\)-tower based at \(x_0\), contradicting the
maximality of \(\Omega_f(x_0)\).

Since \(f\) is progressive, we already have
\[
\omega\leq f(\omega).
\]
The strict inequality \(f(\omega)>\omega\) is impossible, and therefore
\[
f(\omega)=\omega.
\]
Thus \(f\) has a fixed point.
\end{enumerate}
\end{proof}

\end{document}
